\extrachap{Navodila za izdelavo seminarske naloge}

\section*{Tehnična navodila za predlogo}
\begin{itemize}
\item Vaše datoteke so v direktoriju \texttt{SkupinaXX}, kjer \texttt{XX} predstavlja dvomestno številko vaše skupine (npr.\ \texttt{Skupina07}). Glavna datoteka je \texttt{main.tex}.
\item Za začetek uporabite vzorec v direktoriju \texttt{Skupina01}, ki prikazuje zgradbo poglavja, slike, tabele, enačbe, kvantna vezja in programsko kodo.
\item Slike shranite v poddirektorij \texttt{img} svojega direktorija, po možnosti v vektorskem formatu PDF; rastrske slike (npr. posnetke zaslona) v formatu PNG.
\item Vse oznake (\texttt{\textbackslash label}) začnite z \texttt{gXX:}, kjer \texttt{XX} predstavlja številko vaše skupine (npr.\ \texttt{g07:fig:polni\_sestevalnik}).
\item Vire navajajte z ukazom \texttt{\textbackslash cite}. Dodajte jih v skupno datoteko \texttt{references.bib}, ki je v korenskem direktoriju. Če je vir že vsebovan v datoteki, se sklicujte na obstoječi vnos. Ključe novih vnosov začnite z \texttt{gXX:}.
\item Vsak nov vir mora obstajati in biti dosegljiv: navedite polje \texttt{doi} ali \texttt{url}.
\end{itemize}


\section*{Potek dela}

V okviru seminarja se boste ukvarjali s procesnimi platformami, ki za delovanje izkoriščajo kvantne pojave, s kvantnimi celičnimi avtomati (QCA) in s kvantnimi računalniki (QC). Delo poteka v skupinah z največ štirimi študenti. Cilj seminarja so izdelki in poročila, ki jih ob zaključku združimo v skupen zbornik (glej rezultate seminarjev iz prejšnjih let).

\begin{enumerate}
\item \textbf{Prijava teme (rok 15. 11. 2026):} Oddajte enostranski predlog, ki vsebuje naslov, raziskovalno vprašanje, pričakovani rezultat (hipotezo), metrike, s katerimi boste odgovorili na vprašanje, načrt poskusov, potrebna orodja in razdelitev dela med člane.
\item \textbf{Vmesne konzultacije}
\item \textbf{Oddaja (rok 22. 11. 2026):} Oddajte izvorne datoteke poročila v \LaTeX u, poročilo v formatu PDF, prosojnice in povezavo do repozitorija.
\item \textbf{Predstavitev in zagovor (v tednu od 23. 11. 2026:} Vsaka skupina predstavi svoje delo v približno 10 minutah in nato sledijo vprašanja.
\end{enumerate}

Za skupinsko delo uporabljajte repozitorij (npr.\ GitHub ali GitLab). Repozitorij naj vsebuje vse izvorne datoteke (strukture QCA, programsko kodo, skripte za izris grafov) in datoteko \texttt{README} z navodili, kako ponovimo vaše rezultate. Zgodovina sprememb v repozitoriju je del dokazila o vašem delu. Poročilo pišite v okolju \LaTeX, za skupinsko delo lahko uporabite npr. Overleaf\footnote{\url{https://www.overleaf.com/}}. Vire iščite z iskalniki znanstvene literature\footnote{\url{https://scholar.google.com}, \url{https://www.sciencedirect.com}, \url{https://www.scopus.com}, \url{https://arxiv.org}}.


\section*{Zgradba poročila}

Poročilo naj obsega približno 10 strani in ima strukturo znanstvenega članka (glej vzorec v \texttt{Skupina01}):
\begin{itemize}
\item \textbf{Povzetek:} 5 do 10 vrstic z raziskovalnim vprašanjem, metodo in glavnim rezultatom.
\item \textbf{Uvod:} kratek opis področja (največ 2 strani) in jasno zapisano raziskovalno vprašanje.
\item \textbf{Metode:} zasnova poskusov, uporabljena orodja in njihove nastavitve (npr. simulacijski model in parametri v QCADesigner-E, šumni model ali oznaka kvantnega procesorja), metrike.
\item \textbf{Rezultati:} meritve v tabelah in grafih ter primerjava z literaturo, kjer je to mogoče.
\item \textbf{Razprava:} odgovor na raziskovalno vprašanje, omejitve in podpoglavje \emph{Kaj ni delovalo}, v katerem opišete vsaj en neuspešen poskus in razlog zanj.
\item \textbf{Zaključek.}
\item \textbf{Doprinosi avtorjev:} za vsakega člana v dveh do treh stavkih zapišite, kateri del izdelka in poročila je pripravil.
\item \textbf{Izjava o uporabi generativne umetne inteligence:} navedite orodja in namen uporabe (npr. pomoč pri kodi, lektoriranje, iskanje virov).
\end{itemize}
Literatura se izpiše skupaj za ves zbornik na koncu dokumenta.

Uporaba generativne umetne inteligence je dovoljena, vendar ne nadomesti razumevanja. Na zagovoru mora vsak član skupine znati razložiti svoj izdelek in poročilo.


\section*{Izbira teme}

Vsaka tema mora imeti raziskovalno vprašanje, na katero odgovorite z lastnimi poskusi in meritvami. Zgolj opis in implementacija znanega algoritma ali strukture ne zadoščata. Pregled literature (npr. uporaba kvantnih računalnikov na nekem področju) ni samostojna tema.

\section*{Kvantni celični avtomati}
\label{g00:sec:qca}

Tipična tema seminarja QCA je podrobna recenzija ustreznih novejših znanstvenih člankov. V nalogi implementirajte in opišite izbrane strukture. Primer je zasnova in implementacija reverzibilne strukture QCA. Reverzibilno procesiranje je opisano v~\cite{landauer1961, bennett1973, toffoli1980, frank2017}. Primera osnovnih univerzalnih reverzibilnih logičnih operatorjev sta Toffolijeva in Fredkinova vrata. Potreben pogoj za fizično reverzibilnost je logična reverzibilnost: iz logičnih vrednosti izhodov lahko enolično določimo logične vrednosti vhodov, torej med vhodi in izhodi obstaja bijektivna preslikava. Fizično reverzibilna struktura mora izpolnjevati še dodatni pogoj: če spremenimo izhodne celice v vhodne in obratno ter določimo vhodne logične vrednosti (na prej izhodnih celicah), moramo na izhodih (prej vhodnih celicah) dobiti ustrezne vrednosti, določene z bijekcijo.

V seminarjih s področja QCA:
\begin{itemize}
\item Strukture zasnujte in simulirajte z ustreznimi orodji, kot je npr. QCADesigner-E~\cite{walus2004, silltorres2018}\footnote{\url{https://github.com/FSillT/QCADesigner-E}}, ki omogoča tudi izračun disipacije energije. Za avtomatsko snovanje lahko uporabite ogrodje fiction~\cite{walter2019fiction} oziroma orodje MNT Designer\footnote{\url{https://github.com/cda-tum/fiction}, \url{https://github.com/cda-tum/mnt-designer}}, za primerjalna vezja pa zbirko MNT Bench\footnote{\url{https://www.cda.cit.tum.de/mntbench/}}.
\item Za vsako strukturo poročajte število celic, površino, zakasnitev v urinih ciklih, število križanj, disipacijo energije, ... Navedite uporabljeni simulacijski model in njegove parametre.
\item V poglavju Metode opišite, ali ste uporabili ad hoc metodo ali ste formalizirali metodo snovanja, ali prosto določite urino cono vsaki celici ali uporabljate strukturirano urino shemo (npr.\ 2DDWave~\cite{vankamamidi2008} ali USE~\cite{campos2016use}), ...
\item Rezultate primerjajte z vsaj eno objavljeno izvedbo iste funkcije.
\end{itemize}


\newpage
\section*{Kvantni računalniki}
\label{g00:sec:qc}

Kvantno računalništvo za procesiranje izkorišča superpozicijo, interferenco in prepletenost kvantnih bitov~\cite{nielsen2010}. Raziskovalci so predstavili kvantne algoritme (npr.\ Shorov~\cite{shor1997} in Groverjev~\cite{grover1996}), ki imajo manjšo časovno kompleksnost kot najboljši znani klasični algoritmi. Fizično izdelani kvantni računalniki so že rešili nekatere posebej izbrane probleme hitreje od najzmogljivejših klasičnih računalnikov~\cite{arute2019quantum, wu2021strong}. Današnji procesorji so še vedno občutljivi na šum~\cite{preskill2018}. Ključen korak do skalabilnih kvantnih računalnikov je popravljanje napak. Leta 2024 je bila prvič prikazana površinska koda, pri kateri se logična napaka z večanjem kode zmanjšuje~\cite{google2025below}.

V seminarjih s področja QC:
\begin{itemize}
\item Vezja implementirajte v ustreznem okolju, kot je npr. Qiskit~\cite{qiskit2024}\footnote{\url{https://quantum.cloud.ibm.com/docs}} ali v drugem sodobnem okolju (utemeljite izbiro).
\item Rezultate idealnega simulatorja primerjajte z rezultati simulatorja s šumnim modelom in po možnosti z zagonom na realnem kvantnem procesorju. Brezplačni paket IBM Quantum Open Plan\footnote{\url{https://quantum.cloud.ibm.com/docs/en/guides/plans-overview}} omogoča omejen čas na kvantnem procesorju, zato poskuse najprej preizkusite s šumnim modelom. V poročilu navedite oznako procesorja, število ponovitev (angl.\ \emph{shots}) in identifikatorje izvedenih nalog (angl.\ \emph{job ID}).
\item Poročajte o lastnostih vezja pred prevajanjem in po njem (število kubitov, globina, število dvokubitnih vrat).
\item Kjer je mogoče, rezultat primerjajte s klasično rešitvijo istega problema.
\end{itemize}
